% \begin{center}
\thispagestyle{empty}
% {\Large Slovenská technická univerzita v Bratislave}
{\Large \myUniversitySK}
\par\end{center}{\Large \par}

\begin{center}
% {\Large Fakulta informatiky a informačných technológií}
{\Large \myFacultySK}
\par\end{center}{\Large \par}

\smallskip{}

\begin{center}
\myEvidenceNumber
\par\end{center}
\vfill{}

\begin{center}
\textbf{\Large \myName}
\par\end{center}{\Large \par}

\medskip{}


\begin{center}
\textbf{\LARGE \myTitleSK }
\par\end{center}{\huge \par}

\medskip{}


\begin{center}

% {\Large Bakalárska práca / }
%{\Large Bachelor's / Master's Thesis: Conceptual Part}
{\Large Diplomová práca: Konceptuálna časť}
\par\end{center}{\Large \par}

\vfill{}

% Thesis Supervisor: \mySupervisor

% \medskip{}
% \myDate

% \pagenumbering{roman}

% \newpage
% \thispagestyle{empty}
% \mbox{}
% \newpage



\begin{center}
\thispagestyle{empty}
% {\Large Slovenská technická univerzita v Bratislave}
{\Large \myUniversitySK}
\par\end{center}{\Large \par}

\begin{center}
% {\Large Fakulta informatiky a informačných technológií}
{\Large \myFacultySK}
\par\end{center}{\Large \par}

\smallskip{}

\begin{center}
\myEvidenceNumber
\par\end{center}
\vfill{}

\begin{center}
\textbf{\Large \myName}
\par\end{center}{\Large \par}

\medskip{}


\begin{center}
\textbf{\LARGE \myTitleSK }
\par\end{center}{\huge \par}

\medskip{}


\begin{center}

% {\Large Bakalárska práca / }
%{\Large Bachelor's / Master's Thesis: Conceptual Part}
{\Large Diplomová práca: Konceptuálna časť}
\par\end{center}{\Large \par}

\vfill{}

\begin{tabbing}
    \=xxxxxxxxxxxxxxxxxxxxi \=\kill
    \> Študijný program:        \> \myStudyProgramSK  \\
    \> Študijný odbor:          \> \myStudyFieldSK  \\
    \> Miesto vypracovania:     \> FIIT STU v Bratislave, Slovensko \\
    \> Vedúci práce:            \> \mySupervisor    
\end{tabbing}
\myDateSK

% Study programme: \myStudyProgram
% Field of study: \myDegreeCourse
% Training workplace:\\ \myInstitute
% Supervisor: \mySupervisor
% \medskip{}
% \myDate

% Študijný program: \myStudyProgram
% Študijný odbor: \myDegreeCourse
% Miesto vypracovania: \myInstitute
% Vedúci práce: \mySupervisor
% \medskip{}
% \myDate


% \newpage
% \thispagestyle{empty}
% \mbox{}
% \newpage

\newpage
\thispagestyle{empty}

{\center{\textbf{Text zadania:}
Miera využívania umelej inteligencie (UI) vo vzdelávacom procese v posledných rokoch výrazne narastá. Nástroje UI môžu podporovať personalizované vzdelávanie, zvyšovať efektivitu výučby a prispievať k rozvoju digitálnych kompetencií používateľov. Zároveň však prinášajú aj nové etické a pedagogické výzvy.
Analyzujte, akým spôsobom sa v súčasnosti implementujú nástroje umelej inteligencie v školskom prostredí. Realizujte kvantitatívne aj kvalitatívne vyhodnotenie posudzovaných riešení s cieľom identifikovať ich prínosy a obmedzujúce faktory.
Navrhnite riešenie, ktoré umožňuje systematický zber údajov o používaní nástrojov umelej inteligencie v školskom prostredí, kategorizáciu používateľov a meranie zmien vo vybraných schopnostiach používateľov pred a po použití nástrojov umelej inteligencie. Implementujte navrhnuté riešenie a overte ho jeho použitím v školskom prostredí, pričom vyhodnoťte jeho funkčnosť.
Analyzujte zozbierané dáta a vyhodnoťte vplyv používania nástrojov umelej inteligencie na vzdelávací proces. Porovnajte navrhnuté riešenie s dostupnými riešeniami v danej oblasti a formulujte odporúčania pre efektívne a zodpovedné využívanie umelej inteligencie vo vzdelávaní.}}

% \includepdf[pages=-]{assets/zadanie.pdf}